% Section 6: Honest Uncertainties

\section{Honest Uncertainties}
\label{sec:uncertainties}

The preceding sections bound their claims where they arise. This section does
not retract those claims. It collects the remaining uncertainties that matter
to a regulator, architect, or examiner deciding how to use the frame.

\subsection{Scope of the frame}

\textbf{Empty-chair enumeration and weighting remain open.} The six chairs in
the BSA/AML example are illustrative, not exhaustive. Practitioners may identify
different affected parties or assign different weight to their interests. The
frame contributes a discipline---name the absent party, the interest at stake,
and the evidence of representation---but it does not supply a complete roster
or a weighting rule. Disagreement is most useful when those choices are made
explicit and can themselves be examined.

\textbf{Structural and produced absence may not exhaust the possibilities.}
Some absences are mixed, and practice may reveal additional categories. The
distinction is useful where it changes the diagnostic response: a structural
absence may require representation or a proxy, while a produced absence calls
for scrutiny of the design that produced it and the inference drawn from it.
The paper does not claim that every case fits cleanly into one category.

\textbf{Generalization is limited to the contexts examined.} The paper develops
the frame through community-banking examples in BSA/AML monitoring, lending,
and pricing. Application to insurance, healthcare, criminal justice, or social
services is a research proposition, not a result. Even within finance, evidence
that a mechanism appears on one dataset or product does not establish transfer
to another. Sectoral use should begin with new empty-chair enumeration, legal
analysis, and validation of the relevant evidentiary relationships rather than
directly importing the examples here.

\subsection{Conditions of examination}

\textbf{Adversarial and cooperative settings differ.} Examination includes
both: an examiner may test a claim against an institution's incentives, and the
same parties may cooperate to resolve a shared safety or compliance problem.
The paper's strongest warnings about explanation and self-verification concern
settings where incentives are opposed or where a claim must survive challenge.
In cooperative settings, explanation artifacts can support discovery and
communication even when they would not independently verify the proposition at
issue. The method and claimed inference should therefore be evaluated in the
actual examination context.

\textbf{Application requires resources and coordination.} The three diagnostic
questions demand time, expertise, access, and institutional support. Not every
examiner can apply them to every system, and a sophisticated institution can
design formally compliant artifacts that remain difficult to interrogate.
Their practical effect may depend less on exhaustive case-by-case use than on
making the questions predictable enough that institutions design systems able
to answer them. This paper does not establish whether supervisory organizations
can provide the necessary capacity or coordinate a consistent practice.

\textbf{Affected-party access is not uniform.} The ontology gives affected
interests standing in the analysis; it does not imply that every affected party
can or should inspect every artifact. Suspicious-activity confidentiality,
personal data, cybersecurity controls, trade secrets, and the rights of other
customers can limit direct disclosure. Representation may therefore travel
through several contest pathways: a customer-facing explanation and correction
process, protected access for an examiner, independent testing of a proprietary
method, aggregate reporting, or an advocate with appropriate standing. These
pathways are not equivalent. Independent validation can test a bounded
proposition while leaving the affected party unable to identify an error in
their own case; a clear notice can enable challenge while revealing little about
population-level behavior. Applying the frame requires stating who can access
which object, what proposition that access permits them to test, and which
unresolved interest remains represented only indirectly. Whether a particular
access restriction is justified is a legal, security, and proportionality
question outside the ontology itself.

\textbf{Representation is not participation or remedy.} Indirect inspection by
an examiner or independent assessor may represent an affected interest in an
evidentiary inquiry while leaving the affected person without direct voice,
access, or relief. Existing complaint, appeal, correction, and review processes
may permit participation or contest, but this paper neither creates those
processes nor establishes when law should require them. Where none exists, the
ontology can make the absent pathway, the institution able to create it, and the
party bearing its absence legible. That diagnosis is not itself a remedy.

\subsection{Architectural boundaries}

\textbf{The five capabilities are provisional and non-exhaustive.}
Revision-aware temporal capture, evidence binding and exportability,
non-collapsing uncertainty representation, technical-change monitoring, and
institutional-conformance monitoring are generated by this frame's diagnostics.
They are neither a complete governance architecture nor proven necessary for
every deployment. Alternative technical, contractual, independent-assessment,
or compensating-process implementations may answer a diagnostic question more
simply. Implementation may also reveal additional needs.

\textbf{Operational tractability is unproven.} The capabilities may impose storage,
integration, validation, and review costs. Those costs may be material for a
community bank and may interact with confidentiality, retention, cybersecurity,
and vendor constraints. Banks remain responsible for their claims but may lack
leverage to obtain evidence from core-system and AI vendors. The paper has not
built or evaluated a complete implementation, compared alternatives, or shown
that benefits justify costs. The capabilities are design hypotheses with
explicit evidentiary functions, not an implementation mandate.

\textbf{Vendor dependence creates a two-level theory of effect.} A community bank may
identify an evidence gap correctly and still be unable to alter a core platform,
obtain method provenance, or negotiate a bespoke export. Large vendors can
spread implementation cost across institutions, but they also choose product
roadmaps and contractual defaults under incentives different from those of one
bank, examiner, or affected customer. One bank can make the dependency visible,
narrow its claim, evaluate available processes, and record the residual risk.
The collective pathway aggregates those limits: examiners identify unsupported
claims within current authority; regulators and trade groups observe recurrence;
and vendors receive a more consistent capability signal. This paper does not
establish that collective signals will overcome switching costs, concentration,
proprietary constraints, or fragmented regulatory expectations.

Nor does vendor implementation end the inquiry. A standardized export may omit
locally material context; independent validation may be too narrow or too
expensive; a compensating process may shift work to already constrained bank
staff. Confidentiality and cybersecurity can also justify limiting access that
would otherwise improve contestability. The relevant comparison is among
feasible evidence pathways and their residual gaps, not between perfect access
and presumed noncompliance.

\subsection{Who moves first}
\label{sec:uncertainties:first-move}

The largest uncertainty is not whether any participant would find the frame
useful. It is that each participant's first move appears to depend on another's.
A community bank converting silent gaps into written disclosures before an
examination has reason to want evidence that examiners will treat a disclosed
gap with a compensating process as adequate posture rather than as a
self-reported finding. A vendor asked to fund an evidence-export capability has
reason to want evidence that examiners are issuing requests in this shape, since
capability purchases follow examinations that institutions could not answer. An
examiner has authority to test a claimed inference but must be able to name the
statute, regulation, or supervisory expectation under which each evidentiary
object may be compelled, and to state what may be supportably concluded when an
institution answers that the object is vendor-proprietary and unavailable. A
regulator considering guidance must demonstrate workability for a small
institution, which requires evidence about burden that only institutional use
would generate.

Those four dependencies form a cycle. The paper cannot resolve it by asserting
that collective action will aggregate individual practice; that assertion is
what the cycle calls into question.

Two artifacts would break it, and neither requires any participant to act
first. The first is an authority crosswalk: for each evidentiary object this
paper names, the existing authority under which its production may be sought,
and the inference available when it is withheld as proprietary. The second is a
mapping from selected FS~AI~RMF control objectives to the records a conforming
institution would necessarily hold, which would let a party without discovery
request documents by name rather than by description. Both are desk work
against existing law, guidance, and the framework text. Neither is supplied
here. Their absence is the most consequential gap in this paper, and it is a
gap in this paper rather than in the coordination problem it describes.

A preliminary attempt at the first suggests the crosswalk will not divide the
objects into compellable and non-compellable. Supervisory authority is not
confined to inspecting records that happen to exist: reporting powers can
require statements and special reports prepared on demand, recordkeeping
regulation can mandate that records be generated prospectively, and an
examiner's own analysis of produced data is the examiner's work product rather
than something compelled from the institution. The line that does hold is
narrower and concerns time. An institution can be required to answer for a
missing record, to prepare a report, or to build a prospective control. It
cannot be required to produce, after the fact, an object that would exist only
if the decision process had preserved it contemporaneously.

That line cuts through the objects unevenly. Feature attributions, local
approximations, and internal-state observations are in principle recomputable
later. What is not recomputable is the binding: which model version and
configuration, under which policy, actually produced a particular decision on a
particular file. Where that binding was not captured at decision time, the
recomputable objects lose their referent, and the question of whether they
could have been compelled becomes moot --- there is nothing to compel them
against. Decision-time context and the record of what a human reviewer was
actually shown are subject to the same limit and have no recomputable
substitute at all.

The practical consequence is that one narrow record carries most of the weight,
and the precedent for creating it already exists. Mandated loan-level reporting
under the Home Mortgage Disclosure Act did not ask institutions to surrender
data they happened to hold; it required the record to be generated so that
oversight would have something to examine. A comparable and considerably
smaller requirement --- that a decision record bind to the operative model
version, configuration, and policy in force at decision time --- is the
smallest change that would convert several currently unreachable objects into
reachable ones. This paper does not propose that rule, and identifying the
smallest sufficient rule is not the same as establishing that it should be
adopted.

A third artifact would do more and cannot be produced at a desk: one worked
execution of the evidence request of Section~\ref{sec:request} at an actual
institution, recording which fields were answerable from records the bank held,
which required a compensating process, which were unanswerable because a vendor
controlled the evidence, and what the exercise cost in staff hours. That single
exhibit is what a regulator would need for a workability finding and what a
vendor would need for a demand signal.

\subsection{Theory of effect}

The paper assumes that clearer evidentiary distinctions can change what
architects preserve, what institutions claim, and what examiners ask. That
effect is plausible but unproven. Regulatory practice may instead stabilize
around easier documentary proxies; institutions may lack the authority or
resources to change vendor systems; or the vocabulary may be adopted without
the discipline it names. Coordination may produce another dashboard while
leaving underlying evidence inaccessible. Conversely, the frame may be useful even without
formal adoption if it helps one regulator identify an unsupported inference or
one architect preserve evidence that would otherwise disappear.

Adoption can fail by ceremonial compliance as well as rejection. Participants
might add ``affected party,'' ``binding,'' or ``indeterminacy'' fields to a
template while leaving the underlying artifact, access, and decision process
unchanged. Counts of completed fields would then measure use of the vocabulary,
not its effect. Evaluation should instead sample whether a claimed inference can
be traversed to its evidence, whether defeating observations are genuinely
available, whether declared scope reconciles to operational populations, and
whether identified gaps change a decision, disclosure, procurement request, or
policy analysis. Even those observations do not isolate causation without a
comparison, but they are closer to the proposed theory of effect than document
production alone.

The appropriate test is practical and revisable: do the diagnostics reveal
material inference gaps, do the proposed properties make those gaps more
inspectable, and do affected parties receive better representation as a result?
Evidence against any link should narrow or change the frame. This position paper
offers the questions and provisional capabilities; it does not claim the
outcomes of their adoption.
